\chapter{A fresh complete outer-testing certificate for the original R2 source}
\label{component:p8f8780350b13}
\noindent\textbf{Component }\path{p8f8780350b13}\par
{\small\noindent Source: \path{NEW/global/r2_outer_testing_join_REDERIVED.tex}\par}
{\small\noindent SHA256: \path{8f8780350b13c682d2f4e9dc7d20b1dbd0c7b1a3b7a4828d78e472d28838f298}\par}
{\small\noindent Inventory use: R2O1-R2O24\par}
\medskip
\section{A fresh complete outer-testing certificate for the original R2 source}
\label{p8f8780350b13:sec:r2-rederived-outer-join}

This proof treats every independent testing pair outside the positive
double-infinitesimal angle chart of the original R2 source.  It uses an
explicit full-parent Bell reflection and the two-current and six-root
identities that were freshly proved for R1.  Only those universal algebraic
identities and their real matrix arguments are reused.  No assertion that
the reflected source is an original R1 source is made.

\subsection{Actual source and exact reflected parent}

The source is the original weak-C2 tuple, with $d=1$, $\lambda=\beta=0$,
the six coefficients in (RER1), all original strict gaps, $n\geq0$,
$a,B,t_s>0$, and $0<r,h<1$.  R2 means
\[
t_s\text{ unbounded},\qquad B, z=aB, y=t_sB\text{ infinitesimal}.
\tag{R2O1}
\]
Put $S=1/t_s$, $l=1-h$, $e=1-r$, $R_s=1+hz$,
$g_s=R_s-ae>0$, and $k=(1+ar)/(1+a)$.  The rederived exact source
lemma (RER1)--(RER7) gives
\[
k\text{ a positive bounded unit},\qquad k<1,\qquad
SC l\sim SAh\sim D/S\sim k/2,
\qquad \Delta/(1+a)\sim1.
\tag{R2O2}
\]
It also gives $ay$ infinitesimal, $Q\sim1$, and the exact expressions
\[
n=\frac{eQ-BK_s}{2r\Delta},\qquad
m=\frac{a^2e(1+y)c s_c}{2g_s\Delta},\qquad
BK_s=(1+a)y(1+hS)(1+z),
\tag{R2O3}
\]
where $c=1+zr\sim1$ and $s_c=1+B(l+ar)\sim1$.
There is no positive lower bound assumed on $r,h,l$, or $g_s$.

For arbitrary finite testing angles define the unnormalized testing
state transformation
\[
\begin{split}
G\psi=(&uvW-uX-vY+Z,\ -vW+X-uvY+uZ,\\
       &-uW-uvX+Y+vZ,\ W+vX+uY+uvZ).
\end{split}
\tag{R2O4}
\]
It has $G^TG=D_uD_v I$.  The Bell transformation on observables is
\[
(A'_0,A'_1,B'_0,B'_1)=(-B_1,-B_0,-A_1,-A_0).
\]
After the displayed real local basis gauges it sends the angles to
$(u',v')=(v,u)$ and the source coefficients to
\[
(\lambda',n',m',C',A',D',\alpha',\beta')
 =(\alpha,m,n,C,D,A,\lambda,\beta).
\tag{R2O5}
\]
Direct substitution in the whole parent gives
\[
F_{c'}(v,u;G\psi)=D_uD_v F_c(u,v;\psi).
\tag{R2O6}
\]
The checker verifies the coefficient of every original packet in this
single full polynomial identity, as well as $G^TG$.  The transformed
$\lambda'=\alpha$ event is kept as its whole positive square.
Because $D_uD_v>0$ is pure testing, a certificate on the left restores
the right by dividing each root by $\sqrt{D_uD_v}$ within $K_t$.
No source or testing contact is identified in this operation.

In the next two subsections write the transformed testing angles again
as $u,v$, and denote the reflected coefficients by
\[
\mathsf C=C,\quad\mathsf D=A,\quad\mathsf A=D,
\quad\mathsf m=n,\quad\mathsf n=m.
\]
Thus the symbols on the right refer to the original R2 source.
Define the pure-source quantities
\[
g_* =\frac{\mathsf C\mathsf D}{\mathsf C+\mathsf D},\qquad
h_* =\frac{\mathsf D}{\mathsf C+\mathsf D},\qquad
\gamma=\mathsf A g_*.
\tag{R2O7}
\]
If $U_C=SC l$ and $U_A=SAh$, then exactly
\[
Sg_* =\frac{U_CU_A}{hU_C+lU_A}\sim k/2,
\qquad \gamma\sim k^2/4.
\tag{R2O8}
\]
The denominator is a positive bounded unit even when $h$ or $l$ is
infinitesimal.  In particular $1/\mathsf C$, $1/\mathsf D$,
$1/g_*$ are infinitesimal,
$(\mathsf C+\mathsf D)^{-1/2}=O(\sqrt S)$, and $0<h_*<1$.
These are the complete source bounds needed for the first two outer
current identities.

\subsection{Universal outer current identities}

Set
\[
R_t=uW-Y,\quad C_t=v(uX-Z),\quad B_t=X+uZ,\quad
J=R_t+C_t,\quad N=R_t+h_*C_t,
\]
and let $K$ be the exact signed quarter-core
\[
K=\tfrac12J^2+\tfrac12D_u(vY-Z)^2
 +\tfrac12D_vB_t^2-\tfrac12D_uD_vZ^2.
\tag{R2O9}
\]
All unselected original events, including $\lambda'$, are retained
as whole nonnegative packets.

If $1/v$ is bounded, put
\[
K_0=2\sqrt{g_*}C_t,\quad B_0=\sqrt{2D_v/D_u}\,B_t,
\quad c_u=D_u^{-1/2},\quad t_u=uD_u^{-1/2},\quad
\eta_v=\sqrt{D_v}/v.
\]
The entire selected-parent identity is
\begin{align*}
4\mathsf C D_vR_t^2+4\mathsf D J^2+4K
={}&4(\mathsf C+\mathsf D)N^2+4\mathsf C v^2R_t^2
 +2J^2+2D_u(vY-Z)^2\\
&+(K_0,B_0)
\begin{pmatrix}
1-c_u^2\eta_v^2/(2g_*)&t_u\eta_v/\sqrt{2g_*}\\
t_u\eta_v/\sqrt{2g_*}&1
\end{pmatrix}
\binom{K_0}{B_0}.
\tag{R2O10}
\end{align*}
The finite Gram is bounded and has real residue $I$, because $c_u,t_u$
and $\eta_v$ are bounded while $1/g_*$ is infinitesimal.  The actual
separated bounded-Gram lemma makes it tensor SOS.  The roots themselves
need not be bounded.  This gives every state and every other finite angle
when transformed Bob is noninfinitesimal, equivalently when the original
Alice angle is noninfinitesimal.

For positive infinitesimal transformed $v$, put
$q=g_*^{-1/2}$ and define
\[
K_1=2\sqrt{g_*}C_t+qY,\quad
B_1=B_t-uvY,\quad V_0=vY,\quad Y_1=q\sqrt{D_u}Y,
\]
\[
L_0=\tfrac12q c_u(K_1-c_uY_1)-t_uvB_1+c_uvV_0.
\]
Then $L_0=\sqrt{D_u}\,v(vY-Z)$ exactly, and
\begin{align*}
&4\mathsf C D_vR_t^2+4\mathsf D J^2
       +4\mathsf A D_uD_vY^2+4K\\
&\quad=4(\mathsf C+\mathsf D)N^2+4\mathsf C v^2R_t^2
 +2J^2+\mathcal G,\\
\mathcal G={}&K_1^2+2D_v(B_1^2+V_0^2)
 +(4\gamma D_v-c_u^2)Y_1^2\\
&+2qv(K_1-c_uY_1)V_0-2L_0^2.
\tag{R2O11}
\end{align*}
This is the complete four-root Gram, including all normal terms.  Its
residue in the ordered frame $(K_1,B_1,V_0,Y_1)$ is
\[
\operatorname{diag}(1,2,2,k_0^2-c_{u,0}^2).
\tag{R2O12}
\]
Subscript $0$ denotes real residue.  Hence it is strict whenever its
last entry is positive.  This includes every unbounded transformed
$u$ and every bounded positive-unit transformed $u$ when $k_0=1$.
All inverse and square-root factors in these identities belong to one
field at a time.

\subsection{The remaining unit angle and both coefficient openings}

It remains to treat transformed $v$ infinitesimal, transformed $u$ a
bounded positive unit, and $k_0<1$.  The source identity
\[
1-k=\frac{ae}{1+a},\qquad 0<ae<R_s\sim1
\tag{R2O13}
\]
shows that $a,e$ are positive bounded units: if $a$ were infinitesimal
or unbounded, or if $e$ were infinitesimal while $a$ is bounded, then
$1-k$ would be infinitesimal.  Thus in the current chart $a$ is a
positive unit, $e$ is a positive unit, and both $1+a$ and $\Delta$
are units.

Use the opening-safe source coordinates
\[
j=R_s/a,\qquad p=e,\qquad j-p=g_s/a>0.
\tag{R2O14}
\]
They satisfy $0<p<\min(1,j)$, $j,p$ are positive bounded units, and
the residue of $(j+1-p)/(j+1)$ equals $k_0$.  The exact relative
critical coefficient identities are
\[
\frac{\mathsf m}{jp/[2(1-p)(j+1)]}
 =\left(Q-\frac{BK_s}{e}\right)\frac{R_s+a}{R_s\Delta},
\qquad
\frac{\mathsf n}{p/[2(j-p)(j+1)]}
 =(1+y)c s_c\frac{R_s+a}{\Delta}.
\tag{R2O15}
\]
Their right sides have residue one.  Indeed $BK_s/e$ is infinitesimal
here by (R2O3), and all other normalized factors have the limits stated
in (RER).  No division by $r$, $g_s$, $n$, or an infinitesimal gap
was used to assert a bounded error.  The displayed quotients cancel
those openings exactly before the limit is taken.  In particular the
critical chart itself forces original $n>0$, even though the theorem
as a whole retains $n=0$.

Choose positive real caps $m_c\leq\mathsf m$ and
$n_c\leq\mathsf n$, with the choices specified below.  Keep both
nonnegative excesses as the original whole packets.  Define six rows
\[
\widehat N=\sqrt{\mathsf C+\mathsf D}\,N,\quad
P=C_t/\sqrt S,\quad Y_0=\sqrt S\,Y,
\quad X,\quad Z,\quad V=vY.
\tag{R2O16}
\]
They can be unbounded or linearly dependent.  If
$\delta=(\mathsf C+\mathsf D)^{-1/2}$, then
$R_t=\delta\widehat N-h_*\sqrt S P$.  The selected entire
quarter-parent is exactly
\begin{align*}
\mathcal P_{\rm sel}={}&\widehat N^2+\mathsf C v^2R_t^2
 +\tfrac12J^2+(Sg_*)P^2
 +(\mathsf A/S)D_uD_vY_0^2\\
&+m_c\{(v/u)R_t+(D_u/u)V-X-uZ\}^2
 +n_cD_uD_vX^2\\
&+\tfrac12D_u(V-Z)^2+\tfrac12D_v(X+uZ)^2
 -\tfrac12D_uD_vZ^2.
\tag{R2O17}
\end{align*}
The full-state row relation is
\[
PY_0=V(uX-Z).
\tag{R2O18}
\]
Add $2b(PY_0-V(uX-Z))=0$ for a real $b$.  Keep the two displayed
positive squares $\mathsf C v^2R_t^2$ and $J^2/2$ whole.  The
remaining Gram is bounded: $u,1/u$ are bounded, $Sg_*$ and
$\mathsf A/S$ are units, and the $R_t$ contributions in the
$m_c$ square have coefficients $\delta v/u$ and
$h_*\sqrt S v/u$, both infinitesimal.  Its residue is the direct sum
of a unit pivot, the matrix
\[
\begin{pmatrix}k_0/2&b\\b&k_0(1+u_0^2)/2\end{pmatrix},
\tag{R2O19}
\]
and the following real quadratic form on $(X,Z,V)$:
\[
Q_3=m_c\{(D_u/u)V-X-uZ\}^2+n_cD_uX^2
 +\tfrac12D_uV^2-D_uVZ+\tfrac12(X+uZ)^2
 -2bV(uX-Z).
\tag{R2O20}
\]
In this formula $u=u_0>0$ is real.  We now give all finite-cap choices;
no opening of either source coefficient is suppressed.

If $p_0<1$ and $j_0>p_0$, use initially
\[
m_0=\frac{j_0p_0}{2(1-p_0)(j_0+1)},\quad
n_0=\frac{p_0}{2(j_0-p_0)(j_0+1)},\quad b=k_0/2.
\]
The three leading principal minors of $Q_3$ are
\[
m_0+n_0D_u+\tfrac12,\quad
\frac{p_0u_0^2D_u(j_0+1-p_0)}
 {4(j_0-p_0)(j_0+1)^2(1-p_0)},\quad
\frac{p_0u_0^2D_u(j_0+1-p_0)^2}
 {8(j_0-p_0)(j_0+1)^3},
\tag{R2O21}
\]
all strictly positive.  Matrix (R2O19) has determinant
$k_0^2u_0^2/4>0$.  Choose real caps slightly below $m_0,n_0$;
openness of real positive definiteness and (R2O15) ensure they remain
strict and are below the actual coefficients.

If $j_0=p_0<1$, the coefficient $\mathsf n$ is unbounded and
$m_0=j_0^2/[2(1-j_0)(j_0+1)]$ is finite positive.  Take
$b=1/[2(j_0+1)]$.  On the kernel $X=0$ of the $\mathsf n$ packet,
the first diagonal and determinant of $Q_3$ on $(Z,V)$ are
\[
u_0^2(m_0+1/2)>0,\qquad
u_0^2/[4(j_0+1)^2]>0.
\tag{R2O22}
\]
Choose $m_c<m_0$ sufficiently close and then a sufficiently large
real $n_c$.  A real block Schur complement makes the full matrix strict.

If $p_0=1<j_0$, the coefficient $\mathsf m$ is unbounded and
$n_0=1/[2(j_0^2-1)]$ is finite positive.  Choose a real $b$ with
\[
\frac{j_0}{2(j_0+1)}<b<
\min\left\{\frac{k_0\sqrt{D_u}}2,
                 \frac{j_0}{2(j_0-1)}\right\}.
\]
The interval is nonempty because $u_0>0$.  On the kernel
$(D_u/u)V-X-uZ=0$ of the $\mathsf m$ packet, the first diagonal
and determinant in $(X,V)$ are
\[
\frac{D_u}{2(j_0^2-1)}>0,\qquad
-\frac{D_u^2[2b(j_0-1)-j_0][2b(j_0+1)-j_0]}
       {4u_0^2(j_0^2-1)}>0.
\tag{R2O23}
\]
Choose $n_c<n_0$ sufficiently close, then a sufficiently large real
$m_c$.  Matrix (R2O19) and the full $Q_3$ are strict.

If $p_0=j_0=1$, both coefficients are unbounded and $k_0=1/2$.
Choose $1/4<b<\sqrt{D_u}/4$.  On the common event kernel
$X=0$, $Z=D_uV/u_0^2$, the remaining form is
\[
\frac{D_u(4b-1)}{2u_0^2}V^2>0.
\tag{R2O24}
\]
Sufficiently large real caps on both events make the full form strict.
The finite-cap step in each case is the elementary real fact that a
quadratic form positive definite on the kernel of a nonnegative event
Gram becomes positive definite after a sufficiently large multiple of
that event Gram is added; a block Schur complement proves it.

These four cases exhaust $0<p_0\leq\min(1,j_0)$ with positive $j_0,p_0$.
The separated bounded-Gram lemma now proves the entire selected parent
(R2O17), with every state normal retained.  Restore the original excesses
and all unused packets.  The real caps are used only on this testing
chart: no globally physical capped source is asserted, so this argument
does not conflict with the finite-contact cap exclusion (CAP15).

\subsection{Exact outer and projective join}

Return to the original testing angles.  The full-parent sign gauges
(SJ3) permit $u,v\geq0$ without changing source coefficients.  Literal
zeros and projective infinities are the exact signed-core squares
(SJ10), plus the untouched event squares.  For finite positive angles:
\begin{itemize}
\item if the original $u$ is noninfinitesimal, use (R2O10);
\item if the original $u$ is infinitesimal and $v$ is unbounded,
      use (R2O11)--(R2O12);
\item if the original $u$ is infinitesimal and $v$ is a positive
      bounded unit, use (R2O11) when $k_0=1$, and the complete
      six-root argument (R2O13)--(R2O24) when $k_0<1$.
\end{itemize}
These are exactly all finite positive angle pairs outside the case in
which both are infinitesimal.  Every argument is a whole-parent state
Gram, so no additional state chart or state-ratio restriction is present.
Equation (R2O6) restores the original frame and positive pure-testing
factor.  The zero state and the real/imaginary split for complex states
are immediate.  No source-to-testing rate comparison was made.

\paragraph{Assembly interface and remaining obligation.}
The original R2 outer testing region, every state, and all literal
projective endpoints are proved here.  A full-parent certificate on the
remaining original positive $u,v$ double-infinitesimal chart, for the
same unchanged source and all states, combines with this result to give
the whole original R2 row.  That local certificate is a separate
hypothesis until its actual proof is independently accepted.  This
outer theorem does not itself close any other source row or the global
finite semidefinite-lift criterion.
